% Resumo: parágrafo único, no máximo 1160 caracteres (modelo IMT).
Os distúrbios osteomusculares atingem a maioria dos trabalhadores de escritório e estão associados a posturas inadequadas mantidas por longos períodos. Este trabalho desenvolve um sistema embarcado que detecta postura sentada prejudicial e alerta o usuário por vibração. Uma cadeira com oito sensores de força resistivos mede a distribuição do peso; um colar com sensor inercial mede a orientação da parte superior das costas. A cadeira funde as duas fontes, calibrada pela postura ereta de cada usuário, e envia o alerta ao colar por rádio. Em testes com oito participantes, avaliando cada um fora do ajuste dos parâmetros, o sistema alertou em 58 de 62 trechos de postura prejudicial, com alerta em 3\,\% do tempo em postura ereta ou reclinada. A fusão eliminou a maior parte dos falsos alarmes ao reclinar e passou a detectar o escorregamento do quadril. Os resultados indicam a viabilidade técnica da abordagem, com limitações quanto ao tamanho da amostra, à projeção isolada da cabeça e ao consumo de energia.

\vspace{0.2cm}

\textbf{Palavras-chave:} Monitoramento postural. Ergonomia. Sensores de força resistivos. Unidade de medição inercial. Fusão de sensores.
